\section{Discussion}
\label{sec:discussion}

The results show that MCP is a more accurate and reliable solution than RAG, but has the drawback
of additional context usage as well as the potential for extra token usage if caching is not
implemented. This advantage may decrease as more processes are added to the geoserver, as our
example server provided only 25~processes. Paraphrasing did not have a large impact, proving that
these systems can handle a variety of distinct inputs. However, so-called extreme paraphrases
were not tested and all prompts were required to be grammatically correct. Additionally, the
intended workflow had a well-defined correct output and a predictable set of steps to take.
Future work could investigate both extreme paraphrases as well as geospatial analysis that has
open-ended workflows and results.

For this workflow, MCP outperformed RAG assisted tool retrieval by a noticeable margin,
especially when considering the strict workflow criteria. These results persisted across seed
prompt and paraphrased prompt trials. In addition, MCP did not show any drastic differences
between strict and lenient grading, in comparison to RAG which struggled.

On the other hand, the MCP tool list added a significant amount of tokens to the context of any
LLM call. This may be avoided by using token caching with providers that implement this feature.
Then, each repeated token input would be significantly cheaper from an economic point of view.
However, this does not avoid the drawbacks of a larger context usage in general. This includes
degraded performance due to superfluous information in the context window
\citep{anthropic2025b}.

The geoserver in the benchmark had 25~separate processes available as tools for the agent to
call. All processes have been described with standard metadata as required for the OGC API
Processes. Both approaches, MCP and RAG based, used the identical technical set-up with the same
metadata as textual resource to help the LLM select the best suited process. 25~processes is a
reasonable amount to test tool selection, but raises a question about scalability. Would the MCP
performance retain its advantage if there were 50 or more tools available instead? Some research
shows that the limit lies somewhere around 50~tools before the metadata begins to flood the
context window with too much information. Additional metadata, tailored to the LLM selection
process (including e.g.\ counter examples explicit ``Dos and Don'ts'') could further help to
improve the tool selection.

While paraphrases were used to test the variability of prompt inputs, they were semantically
similar in general. They included simple vocabulary swaps, sentence clause reordering, and other
such minor adjustments. It did not include any so-called extreme paraphrases, where the sentence
is changed to such a degree that the underlying meaning may begin to deviate. This presents the
difficulty of ensuring that the same intent is present across prompts but also provides the
opportunity to test the robustness of the models against much more strenuous conditions.

In a similar vein to the degree to which the paraphrases were different from the seed prompt, the
paraphrases were also required to be grammatically correct. This paper chose to focus on
paraphrasing on its own, but the larger aim was to check the agent's performance against a wide
variety of realistic inputs. It is very reasonable to expect grammatical and spelling errors in
prompt input, but generating these would provide a similarly complex challenge as generating
paraphrases. Nonetheless, a test of these types of errors would give additional evidence for or
against the robustness of the agentic workflow generator.

The complexity of the desired task also played a part. The example used in this paper was
relatively simple with a well-defined correct output. There is some but not too much room for
interpretation or variability. This provides a convenient way to test and measure the outputs but
leaves out the potential effects of a more difficult task. As with grammar and extreme
paraphrasing, this is another dimension that may potentially prove more advantageous for RAG over
MCP.
